\chapter{讨论、局限与结论}\label{chap:conclusion}

\section{讨论}\label{sec:conclusion-discussion}

本文提出的 GeneaLink 系统将异构图变换器（HGT）链路预测\citep{hu2020hgt}、子图链路预测（SEAL）的微观图样抽取\citep{zhang2018seal}与基于生成式智能体\citep{park2023generative}的六轮多智能体协商，通过六视图可视分析界面有机串接。其方法论新颖性在于：以"重建交互博弈"取代静态分数排序——传统链路预测仅输出标量并交由用户信任，而本文将拓扑证据转译为夫妻双方智能体可援引的论据，使每一条被接受的婚姻边都对应一段可被审阅、可被反驳的协商轨迹。这一取向亦回应了思维链推理可能流于事后合理化\citep{turpin2023faithfulness}的隐忧：当推理过程不再以模型独白展开，而以双方对抗与对齐落地，史学家便能在其中嵌入领域先验并据此调整。

\section{局限性}\label{sec:conclusion-limit}

\textbf{泛化稳健性。}HGT 编码器在 CMGPD-LN 上训练，其样本密度与编码一致性显著高于多数田野家谱；SEAL 微结构图样集合与宏观协变量（粮价、灾荒等）亦为辽宁特定，跨面板的逐项消融研究尚需在野外完成方能宣称可移植性。

\textbf{大语言模型来源可审计性。}系统将完整智能体对话保存于审计追踪中并在协商竞技场视图 E 中呈现，但这一记录与具体的模型 checkpoint 及 prompt 版本之间不存在密码学绑定，恶意或错配的后端可能产出表面合理却与真实推理调用相偏离的 transcript。

\textbf{单一文化先验。}智能体 persona 隐含八旗内婚与父系延续的清代特定假设；朝鲜族谱与越南家谱所遵循的联姻逻辑迥异，将协议直接迁入将带来系统性偏差，必须重新设计 persona 而非仅替换输入数据。

\textbf{计算开销。}六轮六候选的 Qwen 单次协商在常规硬件上约需 30--45 秒，当队列规模超过 20 时累积等待非平凡，需在批量轮次组织与更轻量模型变体之间折中权衡。

\section{未来工作}\label{sec:conclusion-future}

形成性研究与专家评估暴露的若干路径值得后续探索。其一，persona 多样性扩展：在夫妻两侧智能体之外引入户主与社区族老智能体，为跨户与跨社区边的候选提供第三方视角。其二，密码学审计追踪：为每次 MAS 提交记录 prompt 与模型版本字符串的 SHA-256 哈希，使显示 transcript 与真实推理调用之间事后可验证。其三，跨文化迁移：将图模式与 persona 先验改造以适配朝鲜族谱与越南家谱，检验本文架构是否超越辽宁语料而具普适性。其四，主动学习闭环：每次专家会话后增量微调 HGT，使 recall@1 提升在当前基线之上持续累积。

\section{结论}\label{sec:conclusion-final}

综上所述，GeneaLink 使家谱母家婚姻边重建成为协作、可解释、可审计的工作流程：HGT 链路预测在测试集合上达到 Hit@10$=0.624$，配合基于 SEAL 图样与档案事件记录的六轮协商以及六视图联动界面，专家得以承诺其能解释的边、撤销其后悔的决策。镜像家谱缺失模式的边消融实验进一步证实，作为一手史料中近乎不可见的母家女性，在 CMGPD-LN 关系结构中携带高达 62.6\% recall@1 的可恢复信号；而要将这一信号落到史学叙事，机器推理与人工判读缺一不可——专家研究中 5 名受访者在七分制 Likert 量表上对系统给出 6--6--5 的"有用性—可解释性—对纯 MAS 的信任"评分，正是这一互补性的人在环验证。
